\chapter{रासायनिक संदेशवाहक और संकेत-पारक्रमण}\label{ch:b2:cell-signalling}

ऐड्रिनैलिन का कोई एक अणु, यकृत की किसी कोशिका की सतह पर पहुँचकर,
उसमें कभी नहीं घुसता। वह बाहर किसी प्रोटीन से बँधता है, उस
प्रोटीन का आकार बदल देता है, और सेकंड भर के भीतर वह कोशिका दस
हज़ार अणु \omterm{prop:b2:cell-signalling:gpcr}{द्वितीय संदेशवाहक} बना चुकी होती है, एक लाख
एंजाइम-अणु सक्रिय कर चुकी होती है, और ग्लूकोज़ के दस लाख अणु
छोड़ना शुरू कर चुकी होती है --- यानी दस लाख गुनी प्रवर्धनों की एक शृंखला, जिसे
ऐसे किसी संकेत ने शुरू किया जो द्वार के बाहर ही रहा। यह अध्याय उन
संदेशवाहकों के बारे में है जो कोशिकाएँ एक-दूसरे को भेजती हैं और इस बारे में कि कोई कोशिका अपनी
सतह पर आए किसी संदेश को अपने रसायन, अपने जीनों या अपने
आकार में किसी परिवर्तन में कैसे बदलती है: यानी वे ग्राही, वे \omterm{prop:b2:cell-signalling:gpcr}{द्वितीय संदेशवाहक}, काइनेज़ों के वे
सोपान, और वे दो अभिकल्प --- तंत्रिकाओं का तेज़, तार वाला संकेतन
और \omterm{def:b2:cell-signalling:kinds}{हॉर्मोनों} का धीमा, प्रसारित संकेतन --- जिनका विरोध
इस पुस्तक के शेष भर चलता है।

\section{संदेशवाहक}

\begin{definition}[रासायनिक संकेतन के प्रकार]\label{def:b2:cell-signalling:kinds}
\index{हॉर्मोन}\index{अंतःस्रावी संकेतन}\index{पराक्राइन संकेतन}\index{तंत्रिकासंचारी}
कोई कोशिका किसी दूसरी को किसी स्रवित अणु से संकेत देती है, यानी किसी
\emph{संदेशवाहक} या \emph{लिगैंड} से, जो उस लक्ष्य पर या उसमें किसी \emph{ग्राही} से
बँधता है। परास से: \emph{अंतःस्रावी} संकेतन, जिसमें कोई
\emph{हॉर्मोन} किसी ग्रंथि से \omterm{def:b2:blood-circulation:blood}{रक्त} में छोड़ा जाता है और देह की
हर कोशिका तक पहुँचता है, पर काम केवल उन्हीं पर करता है जिनके पास वह ग्राही हो,
और वह भी मिनटों से घंटों में (इंसुलिन, कॉर्टिसोल, ऐड्रिनैलिन, थायरॉक्सिन);
\emph{पराक्राइन} संकेतन, जिसमें वह संदेशवाहक मिलीमीटर भर के भीतर
पड़ोसी कोशिकाओं तक विसरित होता है और रास्ते में नष्ट हो जाता है
(वृद्धि कारक, हिस्टामिन, प्रोस्टाग्लैंडिन,
\cref{ch:b2:vertebrate-development} के आकारजन); \emph{ऑटोक्राइन}, यानी कोशिका का
स्वयं को संकेत देना; और \emph{सिनैप्टिक}, जिसमें कोई तंत्रिकाकोशिका
\qty{20}{nm} के किसी अंतराल के आरपार किसी एक लक्ष्य कोशिका पर कोई \emph{तंत्रिकासंचारी} छोड़ती है,
और वह भी किसी मिलीसेकंड में (\cref{ch:b2:neurons-synapses})।
रसायन से: \emph{जल-घुलनशील} संदेशवाहक --- ऐड्रिनैलिन जैसे अमीनो-अम्ल
व्युत्पन्न, और इंसुलिन जैसे पेप्टाइड तथा प्रोटीन --- वह
झिल्ली पार नहीं कर सकते और सतही ग्राहियों पर काम करते हैं;
\emph{वसा-घुलनशील} संदेशवाहक --- स्टेरॉइड (कॉर्टिसोल,
ईस्ट्राडायोल, टेस्टोस्टेरॉन), थायरॉक्सिन, रेटिनोइक अम्ल, नाइट्रिक ऑक्साइड ---
उसे पार करके भीतर के ग्राहियों पर काम करते हैं। अंतःस्रावी तथा तंत्रिका
तंत्र वे दो रास्ते हैं जिनसे $10^{13}$ कोशिकाओं की कोई देह स्वयं को
समन्वित करती है: एक सबको धीरे प्रसारित करता है, दूसरा
किसी एक तक जल्दी तार जोड़ता है।
\end{definition}

\begin{omfigure}
\begin{tikzpicture}[every node/.style={font=\scriptsize}, >=stealth]
  % endocrine
  \begin{scope}
    \draw[thick, fill=omProp!25] (0,1.6) circle (0.35); \node[above] at (0,2.0) {ग्रंथि};
    \draw[very thick, omThm!60] (-0.5,0.6) -- (4.5,0.6); \draw[very thick, omThm!60] (-0.5,0.2) -- (4.5,0.2); \node[omThm!70!black] at (2,0.4) {रक्त};
    \foreach \x in {0.6,1.5,2.4,3.3} \fill[omProp] (\x,0.4) circle (0.05);
    \draw[->, omProp] (0,1.25) -- (0,0.65);
    \draw[thick, fill=omDef!20] (3.8,-0.8) circle (0.35); \node[below, align=center] at (3.8,-1.2) {ग्राही वाली\\लक्ष्य कोशिका};
    \draw[->, omProp] (3.8,0.15) -- (3.8,-0.4);
    \draw[thick, fill=gray!20] (1.0,-0.8) circle (0.35); \node[below, align=center] at (1.0,-1.2) {ग्राही नहीं:\\उपेक्षा करती है};
    \node[align=center] at (2.2,-2.4) {अंतःस्रावी: रक्त में हॉर्मोन,\\पूरी देह, मिनटों से घंटों में};
  \end{scope}
  % paracrine
  \begin{scope}[xshift=6.6cm]
    \draw[thick, fill=omProp!25] (0.8,0.6) circle (0.35);
    \foreach \a in {0,45,...,315} \fill[omProp] (0.8,0.6) ++(\a:0.6) circle (0.04);
    \draw[thick, fill=omDef!20] (1.9,0.9) circle (0.3); \draw[thick, fill=omDef!20] (1.7,-0.2) circle (0.3); \draw[thick, fill=omDef!20] (-0.2,0.0) circle (0.3);
    \node[align=center] at (0.8,-2.4) {पराक्राइन: पड़ोसियों तक\\विसरित, मिलीमीटर भर के\\भीतर नष्ट};
  \end{scope}
  % synaptic
  \begin{scope}[xshift=11.6cm]
    \draw[very thick, omDef] (0,1.8) -- (0,0.5); \draw[thick, fill=omDef!25] (0,0.3) circle (0.25);
    \draw[thick, fill=omExo!30] (0,-0.5) circle (0.3);
    \foreach \x in {-0.08,0,0.08} \fill[omProp] (\x,-0.05) circle (0.03);
    \node[align=center] at (0.2,-2.4) {सिनैप्टिक: एक लक्ष्य,\\\qty{20}{nm},\\एक मिलीसेकंड};
  \end{scope}
\end{tikzpicture}

{\small रासायनिक संकेतन के तीन परास: कोई \omterm{def:b2:cell-signalling:kinds}{हॉर्मोन} जो
\omterm{def:b2:blood-circulation:blood}{रक्त} से उस हर कोशिका तक प्रसारित हो जिसके पास उसका ग्राही है; कोई पराक्राइन कारक
जो अपने पड़ोसियों तक पहुँचे; कोई \omterm{def:b2:cell-signalling:kinds}{तंत्रिकासंचारी} जो किसी
सिनैप्स के आरपार एक कोशिका तक पहुँचाया जाए।}
\end{omfigure}

\begin{theorem}[ग्राही-अधिग्रहण]\label{thm:b2:cell-signalling:occupancy}
\index{वियोजन नियतांक}\index{ग्राही-अधिग्रहण}
कोई ग्राही $R$ अपने लिगैंड $L$ को प्रतिवर्ती ढंग से बाँधता है, $R + L
\rightleftharpoons RL$, और वह भी किसी \emph{\omterm{thm:b2:cell-signalling:occupancy}{वियोजन नियतांक}} $K_d =
[R][L]/[RL]$ के साथ। साम्य पर अधिग्रहीत ग्राहियों का अंश है
\[
\theta = \frac{[RL]}{[R]_{\text{कुल}}} = \frac{[L]}{K_d + [L]},
\]
यानी वही अतिपरवलय जो किसी एंजाइम की संतृप्ति का है। $K_d$ वह
सांद्रता है जिस पर आधे ग्राही अधिग्रहीत हैं और जो
\emph{बंधुता} मापती है: जितनी नीची, उतनी कसी। हॉर्मोन-ग्राहियों की
$K_d$ नैनोमोलर से पिकोमोलर परास में है, जो उन
सांद्रताओं से मेल खाती है जिन पर \omterm{def:b2:cell-signalling:kinds}{हॉर्मोन} परिसंचरित होते हैं ($10^{-12}$ से
$10^{-9}$ mol/L): $[L] = K_d/10$ पर दसवाँ भाग अधिग्रहीत है, और
$10K_d$ पर नौ दसवाँ। कोई कोशिका ऐसे किसी \omterm{def:b2:cell-signalling:kinds}{हॉर्मोन} का उत्तर नहीं दे सकती जिसके
ग्राही उसके पास न हों, और वह अपनी संवेदिता ग्राहियों की
संख्या बदलकर सधाती है, या --- जैसा \cref{ch:b2:reproductive-hormones} में ---
उन्हें हटाकर जब वह \omterm{def:b2:cell-signalling:kinds}{हॉर्मोन} लगातार मौजूद हो।
\end{theorem}
\begin{proof}
$[R] = [R]_{\text{कुल}} - [RL]$ के साथ, $K_d = ([R]_{\text{कुल}} -
[RL])[L]/[RL]$, इसलिए $[RL](K_d + [L]) = [R]_{\text{कुल}}[L]$, और वही
वह सूत्र है। वह वर्ष 1 के खंड की माइकेलिस--मेंटेन अभिव्यक्ति है,
बिना उस उत्प्रेरक चरण के।
\end{proof}

\begin{omfigure}
\begin{tikzpicture}
\begin{axis}[
  width=0.68\linewidth, height=5.2cm,
  axis x line=bottom, axis y line=left,
  xmode=log, xmin=0.01, xmax=100, ymin=0, ymax=1.1,
  xlabel={लिगैंड सांद्रता ($K_d$ की इकाइयों में)}, ylabel={अधिग्रहीत ग्राहियों का अंश},
  xlabel style={font=\small, at={(axis description cs:0.5,-0.13)}, anchor=north},
  ylabel style={font=\small},
  tick label style={font=\small}, xtick={0.01,0.1,1,10,100}, ytick={0,0.5,1},
  clip=false,
]
\addplot[very thick, omDef, domain=0.01:100, samples=100] {x/(1+x)};
\draw[dashed, gray] (axis cs:1,0) -- (axis cs:1,0.5) -- (axis cs:0.01,0.5);
\node[font=\scriptsize, anchor=north west] at (axis cs:1.1,0.48) {$[L] = K_d$: आधे अधिग्रहीत};
\node[font=\scriptsize, anchor=south east] at (axis cs:0.1,0.1) {$K_d/10$ पर 10\%};
\node[font=\scriptsize, anchor=north west] at (axis cs:10,0.92) {$10\,K_d$ पर 90\%};
\end{axis}
\end{tikzpicture}

{\small किसी लघुगणकीय पैमाने पर लिगैंड-सांद्रता के सापेक्ष
ग्राही-अधिग्रहण: किसी कोशिका की अनुक्रिया $K_d$ पर केंद्रित सांद्रता के
सौ गुना परास पर दसवें भाग से नौ दसवें तक चलती है।}
\end{omfigure}

\section{सतह पर ग्राही}

\begin{proposition}[G-प्रोटीन-युग्मित ग्राही और चक्रीय AMP]\label{prop:b2:cell-signalling:gpcr}
\index{G प्रोटीन}\index{G-प्रोटीन-युग्मित ग्राही}\index{चक्रीय AMP}\index{द्वितीय संदेशवाहक}\index{प्रोटीन काइनेज़ A}
ग्राहियों का सबसे बड़ा कुल --- किसी मनुष्य में कोई \num{800} जीन,
\omterm{def:b2:cell-signalling:kinds}{हॉर्मोनों}, तंत्रिकासंचारियों, गंधों, प्रकाश के लिए --- ऐसे प्रोटीन हैं जो
झिल्ली सात बार पार करते हैं। बाहर बँधता कोई लिगैंड
भीतर उनका आकार बदल देता है, जहाँ वे कोई \emph{\omterm{prop:b2:cell-signalling:gpcr}{G प्रोटीन}} सक्रिय करते हैं: यानी
कोई तीन-उपइकाई वाला स्विच जो अपना बँधा GDP GTP से बदल लेता है, टूटता है, और
अपने GTP-बँधे रूप में कुछ सेकंडों तक किसी प्रभावक एंजाइम को नियंत्रित करता है
और फिर उस GTP का जल-अपघटन करके स्वयं को बंद कर लेता है। चिरपरिचित
प्रभावक \emph{ऐडिनिलिल साइक्लेज़} है, जो ATP से \emph{\omterm{prop:b2:cell-signalling:gpcr}{चक्रीय AMP}}
बनाता है: यानी कोई छोटा, विसरणशील \emph{\omterm{prop:b2:cell-signalling:gpcr}{द्वितीय संदेशवाहक}} जिसकी
सांद्रता सेकंडों के भीतर \qty{0.1}{\micro mol/L} से चढ़कर
कई माइक्रोमोलर हो जाती है और जो \emph{\omterm{prop:b2:cell-signalling:gpcr}{प्रोटीन काइनेज़ A}} सक्रिय करता है, यानी
ऐसा एंजाइम जो दर्जनों लक्ष्य प्रोटीनों का फ़ॉस्फ़ोरिलन --- और इस तरह उन्हें चालू या बंद ---
करता है। उस संकेत का अंत फ़ॉस्फ़ोडाइएस्टरेज़ करता है, जो
cAMP नष्ट करता है, और फ़ॉस्फ़ेटेज़, जो वे फ़ॉस्फ़ेट हटाते हैं।
यकृत की किसी कोशिका पर ऐड्रिनैलिन ऐसे ही काम करता है: ग्राही $\to$ \omterm{prop:b2:cell-signalling:gpcr}{G प्रोटीन}
$\to$ साइक्लेज़ $\to$ cAMP $\to$ काइनेज़ A $\to$ फ़ॉस्फ़ोरिलेज़ काइनेज़
$\to$ ग्लाइकोजन फ़ॉस्फ़ोरिलेज़ $\to$ ग्लूकोज़। ग्लूकागॉन वही
रास्ता लेता है; दूसरे \omterm{prop:b2:cell-signalling:gpcr}{G प्रोटीन} उस साइक्लेज़ को संदमित करते हैं, या
फ़ॉस्फ़ोलिपेज़ C सक्रिय करते हैं, जो दो और \omterm{prop:b2:cell-signalling:gpcr}{द्वितीय संदेशवाहक} छोड़ता है,
\emph{इनोसिटॉल ट्राइफ़ॉस्फ़ेट} --- जो कैल्सियम-भंडार खोलता है --- और
डाइएसिलग्लिसरॉल।
\end{proposition}
\begin{proof}[प्रमाण-सामग्री]
सदरलैंड (1957) ने दिखाया कि यकृत के किसी समांगीकरण में जोड़ा गया
ऐड्रिनैलिन ग्लूकोज़ केवल तभी छोड़ता था जब झिल्लियाँ मौजूद हों, कि
वे झिल्लियाँ कोई ऊष्मा-स्थायी कारक बनाती थीं जो अकेला ही
झिल्ली-मुक्त प्रभाज में फ़ॉस्फ़ोरिलेज़ सक्रिय कर देता था, और कि वह कारक
\omterm{prop:b2:cell-signalling:gpcr}{चक्रीय AMP} था --- यानी पहला \omterm{prop:b2:cell-signalling:gpcr}{द्वितीय संदेशवाहक}, और यह पहला
प्रमाण कि कोई \omterm{def:b2:cell-signalling:kinds}{हॉर्मोन} किसी अंतःकोशिकीय मध्यस्थ से होकर
कोशिका-सतह पर काम करता है। रॉडबेल और गिलमैन (1970 के दशक) ने पाया कि
उस साइक्लेज़ को GTP चाहिए था और वह \omterm{prop:b2:cell-signalling:gpcr}{G प्रोटीन} अलग किया जो उस
संकेत को ग्राही तथा एंजाइम के बीच ढोता है; और हैज़ा-विष, जो उस
\omterm{prop:b2:cell-signalling:gpcr}{G प्रोटीन} को उसकी चालू दशा में जकड़ देता है, आँत से प्रतिदिन लीटरों तरल
स्रवित कराता है --- यानी वह रोग कोई संकेतन-दोष है।
\end{proof}

\begin{omfigure}
\begin{tikzpicture}[every node/.style={font=\scriptsize}, >=stealth]
  % membrane
  \fill[omExo!30] (0,2.0) rectangle (12,2.5);
  \node[anchor=east] at (-0.1,2.25) {झिल्ली};
  % receptor
  \draw[thick, fill=omDef!40] (1.0,1.7) rectangle (1.8,2.8); \node[anchor=east] at (0.95,2.65) {ग्राही};
  \fill[omProp] (1.4,3.35) circle (0.12); \node[omProp, above] at (1.4,3.45) {ऐड्रिनैलिन};
  \draw[->, omProp] (1.4,3.23) -- (1.4,2.85);
  % G protein
  \draw[thick, fill=omDef!25] (2.6,1.5) ellipse (0.72 and 0.36); \node at (2.6,1.5) {G प्रोटीन};
  \node[below] at (2.6,1.1) {GDP $\to$ GTP};
  \draw[->, thick] (1.8,2.0) -- (2.2,1.75);
  % cyclase
  \draw[thick, fill=omDef!40] (4.2,1.7) rectangle (5.2,2.8); \node[above] at (4.7,2.85) {ऐडिनिलिल साइक्लेज़};
  \draw[->, thick] (3.15,1.5) -- (4.2,1.85);
  % cAMP
  \node[omThm, anchor=west] at (5.35,1.55) {ATP $\to$ cAMP};
  \draw[->, thick, omThm] (4.7,1.7) -- (5.2,1.15);
  % PKA
  \draw[thick, fill=omThm!20] (7.3,0.9) ellipse (0.7 and 0.35); \node at (7.3,0.9) {काइनेज़ A};
  \draw[->, thick, omThm] (6.4,0.9) -- (6.6,0.9);
  % targets
  \draw[->, thick] (8.0,0.9) -- (9.0,0.9);
  \node[anchor=west, align=left] at (9.0,0.9) {फ़ॉस्फ़ोरिलेज़ काइनेज़ $\to$ फ़ॉस्फ़ोरिलेज़\\$\to$ ग्लाइकोजन $\to$ ग्लूकोज़};
  % amplification annotations
  \node[gray, align=center] at (1.4,0.3) {1 हॉर्मोन};
  \node[gray, align=center] at (4.7,0.3) {$\sim 10^{2}$ G प्रोटीन};
  \node[gray, align=center] at (7.3,0.3) {$\sim 10^{4}$ cAMP};
  \node[gray, align=center] at (10.2,0.3) {$\sim 10^{6}$ ग्लूकोज़};
  % off switches
  \node[omExo!70!black, align=center, anchor=south] at (9.6,2.7) {बंद: GTP का जल-अपघटन (सेकंडों में),\\फ़ॉस्फ़ोडाइएस्टरेज़ cAMP नष्ट करता है,\\फ़ॉस्फ़ेटेज़ वे फ़ॉस्फ़ोरिलन उलट देते हैं};
\end{tikzpicture}

{\small ऐड्रिनैलिन से ग्लूकोज़ तक \omterm{prop:b2:cell-signalling:gpcr}{चक्रीय AMP} पथ। हर
चरण अगले की कई प्रतियाँ सक्रिय करता है: \omterm{def:b2:cell-signalling:kinds}{हॉर्मोन} का एक अणु
उत्पाद के दस लाख अणुओं में अंत होता है।}
\end{omfigure}

\begin{theorem}[किसी सोपान में प्रवर्धन]\label{thm:b2:cell-signalling:amplification}
\index{प्रवर्धन (संकेतन)}
यदि किसी सोपान का हर एक $n$ चरण अपने सक्रियण के जीवनकाल भर अगले के
$a$ अणु सक्रिय करे, तो संदेशवाहक का एक अणु
अंतिम उत्पाद के $a^{n}$ अणु बनाता है। सौ-सौ के
चार प्रवर्धक चरणों के साथ --- कोई ग्राही सौ
\omterm{prop:b2:cell-signalling:gpcr}{G प्रोटीन} सक्रिय करता, कोई साइक्लेज़ सौ cAMP बनाता, कोई काइनेज़
सौ एंजाइमों का फ़ॉस्फ़ोरिलन करता, कोई एंजाइम सौ
ग्लूकोज़ अणु छोड़ता --- $a^{n} = 10^{8}$; और व्यवहार में कोई दस लाख भर,
कुछ ही सेकंडों में, $10^{-9}$ mol/L पर किसी \omterm{def:b2:cell-signalling:kinds}{हॉर्मोन} से। प्रवर्धन ही वह है
जिससे \omterm{def:b2:cell-signalling:kinds}{हॉर्मोन} उन उपापचयजों से दस लाख गुना नीची सांद्रताओं पर
असरदार हो सकते हैं जिन्हें वे नियंत्रित करते हैं, और जिससे वह सोपान
हर चरण पर बंद किया जाना चाहिए: ऐसा प्रवर्धक जो बंद न किया जा
सके कोई रोग है (हैज़ा, और वे बहुत-से कैंसर जिनमें कोई
संकेतन प्रोटीन चालू जकड़ा है)।
\end{theorem}
\begin{proof}
चरण 1 $a$ सक्रिय अणु बनाता है; इनमें से हर एक चरण 2 पर $a$
बनाता है, जो $a^{2}$ देता है; और आगमन से, $a^{n}$, यानी $n$ चरणों के बाद।
हर चरण की लब्धि उन क्रियाधार-फेरों की संख्या है जो कोई
सक्रिय उत्प्रेरक निष्क्रिय किए जाने से पहले साध लेता है: यानी उसकी उत्प्रेरक
दर गुणा उसका सक्रिय जीवनकाल।
\end{proof}

\begin{proposition}[ग्राही टाइरोसीन काइनेज़ और कैल्सियम]\label{prop:b2:cell-signalling:rtk}
\index{ग्राही टाइरोसीन काइनेज़}\index{इंसुलिन ग्राही}\index{कैल्सियम संकेतन}\index{कैल्मोडुलिन}
सतही ग्राहियों का दूसरा कुल --- इंसुलिन के लिए, कोशिका-विभाजन चलाने वाले
वृद्धि कारकों के लिए, परिवर्धन के प्रेरक संकेतों के लिए ---
\emph{\omterm{prop:b2:cell-signalling:rtk}{ग्राही टाइरोसीन काइनेज़}} हैं: यानी एक-बार-पार करने वाले
प्रोटीन जो, जब उनका लिगैंड उनमें से दो को साथ ले आए, एक-दूसरे का
टाइरोसीनों पर फ़ॉस्फ़ोरिलन करते हैं और इस तरह किसी
अनुकूलक प्रोटीन-समुच्चय के लिए लंगर-स्थल बना देते हैं। इनसे एक साथ कई सोपान
चलते हैं: तीन काइनेज़ों की एक शृंखला (MAP काइनेज़ सोपान) जो
केंद्रक में अंत होती है और वृद्धि के जीन चालू करती है; कोई लिपिड काइनेज़
जिसका उत्पाद वह काइनेज़ बुलाता है जो इंसुलिन के उपापचयी
प्रभाव मध्यस्थ करता है --- ग्लूकोज़-वाहक झिल्ली तक ले जाना, ग्लाइकोजन-संश्लेषण
चालू करना; और दूसरे। तीसरा सार्वभौम \omterm{prop:b2:cell-signalling:gpcr}{द्वितीय
संदेशवाहक} \emph{कैल्सियम} है: कोशिकाद्रव में
\qty{0.1}{\micro mol/L} पर थामा, यानी बाहर से दस हज़ार गुना नीचे, वह
किसी संकेत या किसी वोल्टता से खुले वाहिकाओं से भीतर बाढ़ की तरह आता है, या
इनोसिटॉल ट्राइफ़ॉस्फ़ेट के उत्तर में अंतर्द्रव्यी जालिका से बाहर,
और मिलीसेकंडों में माइक्रोमोलर तक चढ़ जाता है; \emph{कैल्मोडुलिन} से बँधकर वह
काइनेज़ तथा पंप सक्रिय करता है, और ट्रोपोनिन से बँधकर वह
संकुचन छेड़ता है (\cref{ch:b2:muscle-movement})। पंप उसे सेकंड भर के भीतर उन
भंडारों में लौटा देते हैं। ग्राही जो भी हो, कोशिका का भीतरी भाग
कुछ ही भाषाएँ बोलता है --- cAMP, कैल्सियम, फ़ॉस्फ़ोरिलन --- और वह
विशिष्टता इसमें है कि हर कोशिका-प्रकार के पास उनके काम करने के लिए कौन-से
प्रोटीन हैं।
\end{proposition}

\begin{omfigure}
\includegraphics[width=0.6\linewidth]{images/book4/ai/b4-19-adrenaline-runner.jpg}

{\small वह सोपान काम पर। कोई डर अधिवृक्क मध्यांश से
ऐड्रिनैलिन छोड़ता है; सेकंडों के भीतर \omterm{def:b2:heart:anatomy}{हृदय} दौड़ता है, यकृत
ग्लूकोज़ उड़ेलता है और वायुमार्ग चौड़े होते हैं --- यानी एक संदेशवाहक,
कई ग्राही, और हर कोशिका में दस लाख गुना प्रवर्धन।}
\end{omfigure}

\section{भीतर के ग्राही}

\begin{proposition}[केंद्रकीय ग्राही]\label{prop:b2:cell-signalling:nuclear}
\index{केंद्रकीय ग्राही}\index{स्टेरॉइड हॉर्मोन}
स्टेरॉइड, थायरॉक्सिन, रेटिनोइक अम्ल तथा विटामिन D उस
झिल्ली से विसरित होकर कोशिकाद्रव या केंद्रक में ऐसे ग्राहियों से बँधते हैं जो
स्वयं \emph{\omterm{prop:b2:cell-differentiation:transcriptional}{अनुलेखन कारक}} हैं: वह हॉर्मोन--ग्राही
संकुल विशिष्ट DNA अनुक्रमों से बँधता है और अपने लक्ष्य जीन
चालू या बंद करता है। वह अनुक्रिया धीमी है --- उस mRNA तथा उस प्रोटीन को बनाने में
मिनटों से घंटों --- और टिकाऊ: कॉर्टिसोल घंटों तक
ग्लूकोनवजनन के एंजाइम प्रेरित करता है, ईस्ट्राडायोल दिनों भर गर्भाशय-अस्तर
बनाता है, टेस्टोस्टेरॉन सप्ताहों भर पेशी बनाता है, और थायरॉक्सिन हर कोशिका की
उपापचय दर बाँधता है। कुछ स्टेरॉइडों की सतही ग्राहियों से होकर तेज़
क्रियाएँ भी हैं, और कुछ सतही पथ केंद्रक तक भी पहुँचते हैं
(MAP काइनेज़ सोपान), इसलिए वह भेद निरपेक्ष
नहीं है; पर वह नियम टिका है कि कोई वसा-घुलनशील संदेशवाहक
किसी कोशिका के जीन बदलकर बदल देता है कि वह \emph{है} क्या, जबकि कोई
जल-घुलनशील उसके एंजाइम बदलकर बदल देता है कि वह \emph{करती} क्या है।
\end{proposition}

\begin{omfigure}
\begin{tikzpicture}[every node/.style={font=\scriptsize}, >=stealth]
  \foreach \x/\t in {0/{जल-घुलनशील: ऐड्रिनैलिन, इंसुलिन, ग्लूकागॉन}, 7.4/{वसा-घुलनशील: कॉर्टिसोल, ईस्ट्राडायोल, थायरॉक्सिन}} \node[anchor=south, font=\scriptsize\bfseries] at (\x+2.6,3.4) {\t};
  % left: surface receptor
  \fill[omExo!30] (0,2.4) rectangle (5.2,2.8);
  \draw[thick, fill=omDef!40] (1.6,2.1) rectangle (2.3,3.1);
  \fill[omProp] (1.95,3.35) circle (0.1);
  \draw[->, thick] (1.95,2.1) -- (1.95,1.6) node[below, align=center] {द्वितीय संदेशवाहक,\\काइनेज़: एंजाइम\\सेकंडों में स्विच किए};
  \node[align=center] at (4.5,1.0) {बदलता है कि वह\\कोशिका करती क्या है};
  % right: nuclear receptor
  \begin{scope}[xshift=7.4cm]
    \fill[omExo!30] (0,2.4) rectangle (5.2,2.8);
    \fill[omProp] (1.5,3.35) circle (0.1);
    \draw[->, omProp] (1.5,3.25) -- (1.5,1.9);
    \draw[thick, fill=omThm!20] (1.5,1.5) ellipse (0.72 and 0.3); \node at (1.5,1.5) {ग्राही};
    \draw[thick] (0.6,0.4) -- (4.6,0.4); \draw[thick] (0.6,0.25) -- (4.6,0.25); \node[below] at (2.6,0.2) {DNA: लक्ष्य जीन चालू या बंद, घंटों में};
    \draw[->, thick] (1.5,1.2) -- (2.4,0.5);
    \node[align=center] at (4.0,1.4) {बदलता है कि वह\\कोशिका है क्या};
  \end{scope}
\end{tikzpicture}

{\small दो अभिकल्प। कोई जल-घुलनशील संदेशवाहक सतह पर काम करता है
और सेकंडों के भीतर एंजाइम-क्रियाशीलता बदल देता है; कोई वसा-घुलनशील
भीतर घुसता है, ऐसे किसी ग्राही से बँधता है जो कोई \omterm{prop:b2:cell-differentiation:transcriptional}{अनुलेखन कारक} है, और घंटों भर
जीन-अभिव्यक्ति बदलता है।}
\end{omfigure}

\section{एक साधा हुआ उदाहरण: इंसुलिन, ग्लूकागॉन और ग्लूकोज़ का नियत बिंदु}

\begin{proposition}[दो हॉर्मोन और एक नियत बिंदु]\label{prop:b2:cell-signalling:glucose}
\index{इंसुलिन}\index{ग्लूकागॉन}\index{रक्त ग्लूकोज़}
\omterm{def:b2:blood-circulation:blood}{रक्त} की ग्लूकोज़ \qty{5}{mmol/L} (\qty{0.9}{g/L}) के पास
अग्न्याशयी द्वीपिकाओं के दो विरोधी \omterm{def:b2:cell-signalling:kinds}{हॉर्मोन} थामे रहते हैं। किसी भोजन के बाद
चढ़ती ग्लूकोज़ को $\beta$ कोशिकाएँ भाँपती हैं --- जो
उसका उपापचय करती हैं, अपनी ATP चढ़ते ही कोई पोटैशियम वाहिका बंद कर देती हैं,
विध्रुवित होती हैं, कैल्सियम भीतर लेती हैं और \emph{इंसुलिन} छोड़ती हैं --- और इंसुलिन,
अपने टाइरोसीन-काइनेज़ ग्राही से होकर, पेशी तथा वसा से ग्लूकोज़
ग्रहण कराता है (उनकी झिल्लियों तक वाहक ले जाकर), यकृत से
उसे ग्लाइकोजन तथा वसा के रूप में संचित कराता है, और सारे ऊतकों से वसा जलाना बंद कराता है।
भोजनों के बीच गिरती ग्लूकोज़ $\alpha$ कोशिकाओं को
\emph{ग्लूकागॉन} छोड़ने देती है, जो, यकृत में cAMP से होकर, ग्लाइकोजन
तोड़ता है और अमीनो अम्लों से नई ग्लूकोज़ बनाता है; आपात में ऐड्रिनैलिन वही
तेज़ करता है, और कॉर्टिसोल तथा वृद्धि \omterm{def:b2:cell-signalling:kinds}{हॉर्मोन} घंटों भर।
हर \omterm{def:b2:cell-signalling:kinds}{हॉर्मोन} उस विचलन के अनुपात में छोड़ा जाता है जिसे वह
सुधारता है: यानी दो प्रभावकों वाली कोई ऋण प्रतिपुष्टि, हर दिशा के लिए एक,
और कोई नियत बिंदु जिसकी रक्षा पाँचवें भाग के भीतर तक होती है। मस्तिष्क,
जो प्रतिदिन \qty{120}{g} ग्लूकोज़ जलाता है और कुछ भी संचित नहीं रखता, वही
अंग है जिसकी वह पाश रक्षा करता है; टाइप 1 मधुमेह उन
$\beta$ कोशिकाओं की हानि है, और उसका बिना इलाज वाला मार्ग ---
\qty{20}{mmol/L} पर ग्लूकोज़, इंसुलिन के अभाव में जलती वसा, \omterm{def:b2:blood-circulation:blood}{रक्त} में अम्ल
--- वही है जैसा वह पाश एक भुजा हटा दिए जाने पर दिखता है।
\end{proposition}

\begin{omfigure}
\begin{tikzpicture}
\begin{axis}[
  width=0.82\linewidth, height=5.6cm,
  axis x line=bottom, axis y line=left,
  xmin=0, xmax=6, ymin=0, ymax=1.1,
  xlabel={किसी भोजन के बाद घंटे}, ylabel={शिखर के सापेक्ष},
  xlabel style={font=\small, at={(axis description cs:0.5,-0.12)}, anchor=north},
  ylabel style={font=\small},
  tick label style={font=\small}, xtick={0,1,2,3,4,5,6}, ytick={0,0.5,1},
  legend style={font=\scriptsize, at={(0.5,-0.28)}, anchor=north, draw=none, legend columns=3},
  clip=false,
]
\addplot[very thick, omDef, domain=0:6, samples=80] {0.55 + 0.45*exp(-((x-0.9)^2)/0.5) - 0.08*exp(-((x-4)^2)/2)};
\addlegendentry{रक्त ग्लूकोज़ (0.55 पर 5 mmol/L)}
\addplot[very thick, omThm, domain=0:6, samples=80] {0.15 + 0.85*exp(-((x-1.1)^2)/0.45)};
\addlegendentry{इंसुलिन}
\addplot[very thick, omProp, domain=0:6, samples=80] {0.75 - 0.5*exp(-((x-1.2)^2)/0.8) + 0.25*(1-exp(-x/5))*0.6};
\addlegendentry{ग्लूकागॉन}
\end{axis}
\end{tikzpicture}

{\small किसी भोजन के बाद ग्लूकोज़, इंसुलिन और ग्लूकागॉन। इंसुलिन
उस ग्लूकोज़ के साथ चढ़ता है और उसे नीचे लाता है; और ज्यों ही ग्लूकोज़ उस नियत
बिंदु की ओर गिरती है, त्यों ही ग्लूकागॉन उसे वहीं थामने को चढ़ जाता है।}
\end{omfigure}

\begin{example}[दोनों तंत्रों की गति, पहुँच और लागत]\label{ex:b2:cell-signalling:compare}
\omterm{def:b2:blood-circulation:blood}{रक्त} में ऐड्रिनैलिन लगभग मिनट भर में हर कोशिका तक पहुँचता है
(यानी परिसंचरण-काल) और मिनटों तक काम करता है; कोई तंत्रिका आवेग
मिलीसेकंडों में अपने लक्ष्य तक पहुँचता है और मिलीसेकंडों तक काम करता है। किसी
\omterm{def:b2:cell-signalling:kinds}{हॉर्मोन} का प्रति पहुँची कोशिका दाम लगभग कुछ नहीं --- देह भर फैला कोई नैनोमोल ---
पर वह किसी एक कोशिका को संबोधित नहीं कर सकता; कोई तंत्रिका ठीक एक
कोशिका को संबोधित करती है पर उसे उस तक कोई तार बनाना तथा रखना पड़ता है। देह
दोनों काम में लेती है: \cref{ch:b2:blood-pressure} की अनुकंपी तंत्रिकाएँ सेकंड भर में
\omterm{def:b2:heart:anatomy}{हृदय} तथा \omterm{def:b2:blood-circulation:vessels}{धमनिकाओं} पर काम करती हैं, और अधिवृक्क मध्यांश,
जो स्वयं कोई रूपांतरित अनुकंपी गुच्छिका है, पूरी देह के लिए वही अणु
\omterm{def:b2:blood-circulation:blood}{रक्त} में उड़ेलकर उनका साथ देता है। और सिनैप्स पर
वे दोनों अभिकल्प मिलते हैं, क्योंकि कोई \omterm{def:b2:cell-signalling:kinds}{तंत्रिकासंचारी} वैसा ही संदेशवाहक है जो
उन्हीं क्रियाविधियों से किसी ग्राही पर काम करता है --- कोई वाहिका जो खुले, कोई \omterm{prop:b2:cell-signalling:gpcr}{G
प्रोटीन} जो स्विच करे --- बस किसी मीटर के बजाय बीस नैनोमीटर के
आरपार।
\end{example}

\section{अभ्यास}

\begin{exercise}[$\star$]\label{exo:b2:cell-signalling:1}
अंतःस्रावी, पराक्राइन, ऑटोक्राइन या सिनैप्टिक के रूप में, और
जल- या वसा-घुलनशील के रूप में वर्गीकृत कीजिए: इंसुलिन, कॉर्टिसोल, किसी
सिनैप्स पर ऐसिटिलकोलीन, किसी घाव पर हिस्टामिन, ईस्ट्राडायोल, अधिवृक्क से
ऐड्रिनैलिन, अंतःस्तर से नाइट्रिक ऑक्साइड।
\end{exercise}

\begin{exercise}[$\star$]\label{exo:b2:cell-signalling:2}
यकृत-कोशिका की झिल्ली पर ऐड्रिनैलिन से \omterm{def:b2:blood-circulation:blood}{रक्त} में ग्लूकोज़ तक के चरण
गिनाइए, और हर चरण पर वह बंद-स्विच बताइए।
\end{exercise}

\begin{exercise}[$\star$]\label{exo:b2:cell-signalling:3}
सतही और केंद्रकीय ग्राहियों की तुलना कीजिए: उन्हें क्या बाँधता है, वे कहाँ
हैं, वे कितने तेज़ काम करते हैं, वे क्या बदलते हैं।
\end{exercise}

\begin{exercise}[$\star$]\label{exo:b2:cell-signalling:4}
इस अध्याय में बताए गए तीन सार्वभौम \omterm{prop:b2:cell-signalling:gpcr}{द्वितीय संदेशवाहक} कौन-से हैं,
हर एक कैसे चढ़ाया जाता है, और हर एक कैसे हटाया जाता है?
\end{exercise}

\begin{exercise}[$\star\star$]\label{exo:b2:cell-signalling:5}
किसी ग्राही की $K_d = \qty{2}{nmol/L}$ है। 0.2, 2, 20 तथा
\qty{200}{nmol/L} पर अधिग्रहण निकालिए। कोई कोशिका अपनी ग्राही-संख्या दुगुनी कर लेती है:
क्या बदलता है, वह अधिग्रहण या अधिग्रहीत ग्राहियों की संख्या?
\end{exercise}

\begin{exercise}[$\star\star$]\label{exo:b2:cell-signalling:6}
किसी सोपान के चरणों की लब्धियाँ 50, 200, 100 तथा 300 हैं। कुल
प्रवर्धन निकालिए। यदि वह \omterm{def:b2:cell-signalling:kinds}{हॉर्मोन} \qty{1}{nmol/L} पर हो और
ग्राहियों (प्रति कोशिका $10^{4}$) का \qty{1}{\%} अधिग्रहीत हो, तो
वह कोशिका प्रति सेकंड कितने उत्पाद-अणु बनाती है यदि हर चरण की
लब्धि प्रति सेकंड हो?
\end{exercise}

\begin{exercise}[$\star\star$]\label{exo:b2:cell-signalling:7}
हैज़ा-विष आँत की कोशिकाओं में उत्तेजक \omterm{prop:b2:cell-signalling:gpcr}{G प्रोटीन} को उसकी GTP-बँधी
दशा में जकड़ देता है। चरण-दर-चरण समझाइए कि वह रोगी प्रतिदिन लीटरों
तरल क्यों खोता है, और वे जीवाणु चले जाने के बाद भी वह प्रभाव दिनों तक क्यों
बना रहता है।
\end{exercise}

\begin{exercise}[$\star\star$]\label{exo:b2:cell-signalling:8}
जब कोई वाहिका खुलती है तब \qty{2000}{\micro m^3} की किसी कोशिका में
कोशिकाद्रवी कैल्सियम \qty{0.1}{} से \qty{1}{\micro
mol/L} तक चढ़ जाती है। कितने
कैल्सियम आयन घुसे? यदि कोई एक खुली वाहिका प्रति सेकंड $10^{6}$
आयन गुज़ारे, तो एक वाहिका को कितना समय लगा, और कोई कोशिका
एक साथ बहुत-सी क्यों खोलती है?
\end{exercise}

\begin{exercise}[$\star\star$]\label{exo:b2:cell-signalling:9}
इंसुलिन और ग्लूकागॉन दोनों यकृत पर काम करते हैं, एक किसी टाइरोसीन
काइनेज़ से होकर और एक cAMP से होकर। समझाइए कि यकृत-कोशिका दोनों संकेत कैसे
समाकलित करती है, दोनों ऊँचे हों तब ग्लाइकोजन का क्या होता है, और
उनका अनुपात किसी एक के स्तर से अधिक क्यों मायने रखता है।
\end{exercise}

\begin{exercise}[$\star\star\star$]\label{exo:b2:cell-signalling:10}
ऐड्रिनैलिन यकृत की किसी कोशिका में \qty{5}{s} के भीतर cAMP चढ़ा देता है और
\qty{30}{s} के भीतर ग्लूकोज़-निर्गम; जबकि कॉर्टिसोल यकृत के ग्लूकोनवजनक
एंजाइम \qty{4}{h} भर चढ़ाता है। विसरण, अनुलेखन (किसी
\qty{3000}{}-न्यूक्लिओटाइड जीन पर प्रति सेकंड \qty{50}{}
न्यूक्लिओटाइड), अनुवादन (किसी \num{500}-अवशेष प्रोटीन पर प्रति सेकंड
\qty{5}{} अमीनो अम्ल) और हज़ारों एंजाइम-अणु जमा करने की ज़रूरत से
उन दो काल-मानों का लेखा दीजिए।
\end{exercise}

\begin{exercise}[$\star\star\star$]\label{exo:b2:cell-signalling:11}
कोई \omterm{def:b2:genome-diversification:mutation}{उत्परिवर्तन} किसी वृद्धि-कारक \omterm{prop:b2:cell-signalling:rtk}{ग्राही टाइरोसीन काइनेज़} को
अपने लिगैंड के बिना ही द्विलक होने, और इस तरह संकेत देने, देता है। उस कोशिका के व्यवहार की
भविष्यवाणी कीजिए, समझाइए कि यह कैंसर में कोई आम चरण क्यों है, और कहिए
कि उस काइनेज़ के ATP स्थल को रोकने वाली कोई दवा क्या करती।
\end{exercise}

\begin{exercise}[$\star\star\star$]\label{exo:b2:cell-signalling:12}
``अंतःस्रावी तंत्र प्रसारित करता है; तंत्रिका तंत्र तार जोड़ता है।''
दोनों अभिकल्पों पर गति, पहुँच, विशिष्टता तथा लागत में विचार कीजिए, और
कहिए कि हर एक कहाँ अपरिहार्य है और देह उसी संदेश के लिए दोनों
कहाँ काम में लेती है।
\end{exercise}

\section{समस्या: रक्त से जीन तक एक हॉर्मोन}

\begin{problem}[{सप्तांत समस्या --- ऐड्रिनैलिन अधिवृक्क ग्रंथि से यकृत
की किसी कोशिका की ग्लूकोज़ तक पीछा किया गया, और उसकी सांद्रता,
ग्राही-अधिग्रहण, सोपान-लब्धि तथा समय गिने गए, फिर कॉर्टिसोल
किसी जीन तक पीछा किया गया, और ग्लूकोज़-पाश का विश्लेषण, और अंत उस
अधिग्रहण, उस प्रवर्धन तथा हर हॉर्मोन के लिए गए समय पर}]
\label{pb:b2:cell-signalling:1}
आँकड़े: किसी डर में अधिवृक्क मध्यांश \qty{5}{L} \omterm{def:b2:blood-circulation:blood}{रक्त} में
\qty{1}{nmol} ऐड्रिनैलिन छोड़ता है; यकृत-ग्राही पर ऐड्रिनैलिन की $K_d$
\qty{1}{nmol/L} है; यकृत की कोई कोशिका (\qty{5000}{\micro m^3})
$2\times 10^{4}$ ग्राही रखती है। सक्रियण की प्रति सेकंड सोपान-लब्धियाँ:
ग्राही $\to$ \omterm{prop:b2:cell-signalling:gpcr}{G प्रोटीन} 20, \omterm{prop:b2:cell-signalling:gpcr}{G प्रोटीन} $\to$ cAMP 100
(यानी \omterm{prop:b2:cell-signalling:gpcr}{G प्रोटीन} के 3-सेकंड जीवनकाल में उस साइक्लेज़ का फेर),
cAMP $\to$ काइनेज़ A 1 (रससमीकरणमितीय), काइनेज़ A $\to$ फ़ॉस्फ़ोरिलेज़
काइनेज़ 50, $\to$ फ़ॉस्फ़ोरिलेज़ 50, $\to$ ग्लूकोज़ 1000। यकृत
\qty{100}{g} ग्लाइकोजन थामे है। कॉर्टिसोल: \qty{4000}{}
न्यूक्लिओटाइड के किसी जीन पर प्रति सेकंड \qty{50}{} न्यूक्लिओटाइड पर
अनुलेखन, किसी \num{600}-अवशेष एंजाइम पर प्रति सेकंड \qty{5}{} अवशेष पर
अनुवादन, प्रति mRNA \num{20} राइबोसोम, प्रति घंटा \num{100}
mRNA बने, और ऐसा प्रभाव जिसे $10^{5}$ एंजाइम-अणु चाहिए।

\medskip
\noindent\textbf{भाग I --- वह \omterm{def:b2:cell-signalling:kinds}{हॉर्मोन} पहुँचता है।}
\begin{enumerate}
  \item उस छोड़े जाने के बाद ऐड्रिनैलिन की प्लाज़्मा-सांद्रता
        निकालिए, उसका हटना छोड़ते हुए।
  \item ग्राही-अधिग्रहण और प्रति यकृत-कोशिका अधिग्रहीत
        ग्राहियों की संख्या निकालिए।
  \item ऐड्रिनैलिन \qty{2}{min} के अर्ध-आयु से हटाया जाता है।
        \qty{6}{min} बाद उसकी सांद्रता तथा अधिग्रहण क्या है?
  \item लगभग मिनट भर का परिसंचरण-काल देखते हुए ऐड्रिनैलिन को
        अधिवृक्क से यकृत तक पहुँचने में कितना समय लगता है? यकृत की
        अनुकंपी तंत्रिका उससे तेज़ क्यों काम करती है?
  \item $2\times 10^{5}$ ग्राहियों तथा उसी $K_d$ वाली कोई कोशिका:
        क्या बदलता है?
  \item $K_d$ परिसंचरित सांद्रता से मेल क्यों खिलाई जाती है ---
        $K_d = \qty{1}{\micro mol/L}$ वाला कोई ग्राही किसी
        नैनोमोलर \omterm{def:b2:cell-signalling:kinds}{हॉर्मोन} का क्या करता?
\end{enumerate}

\medskip
\noindent\textbf{भाग II --- वह सोपान।}
\begin{enumerate}[resume]
  \item एक अधिग्रहीत ग्राही से प्रति सेकंड ग्लूकोज़-अणुओं तक
        कुल लब्धि निकालिए।
  \item पहले सेकंड में वह कोशिका जितने cAMP अणु बनाती है, और
        उसके आयतन में वे जिस सांद्रता तक पहुँचते हैं, निकालिए।
  \item प्रश्न 2 के अधिग्रहण पर यकृत की एक कोशिका का प्रति सेकंड,
        और प्रति मिनट, ग्लूकोज़-निर्गम निकालिए।
  \item यकृत में $2\times 10^{11}$ कोशिकाएँ हैं। वे लब्धियाँ जो
        निर्गम बताती हैं उसे ग्राम प्रति मिनट में निकालिए, और वास्तव में देखे गए
        \qty{5}{g/min} से तुलना कीजिए। असली निर्गम को क्या सीमित करता है?
  \item इस दर पर \qty{100}{g} ग्लाइकोजन कितनी देर चलता? वह किसी डर में
        चुक क्यों नहीं जाता?
  \item हर प्रवर्धक चरण पर वह बंद-स्विच बताइए और कहिए कि कौन
        सबसे तेज़ काम करता है।
  \item कोई फ़ॉस्फ़ोडाइएस्टरेज़ संदमक (कैफ़ीन) मौजूद है। उस सोपान पर
        और उस ग्लूकोज़-निर्गम पर प्रभाव की भविष्यवाणी कीजिए।
\end{enumerate}

\medskip
\noindent\textbf{भाग III --- कॉर्टिसोल से किसी जीन तक।}
\begin{enumerate}[resume]
  \item एक mRNA का अनुलेखन करने का काल निकालिए।
  \item एक एंजाइम-अणु का अनुवादन करने का काल, और प्रति mRNA प्रति घंटा
        एंजाइम-अणुओं की संख्या निकालिए (हर mRNA घंटे भर
        जीता, और राइबोसोम एक के बाद एक शुरू होते)।
  \item उन \num{100} mRNA से प्रति घंटा बने एंजाइम-अणु, और
        $10^{5}$ तक पहुँचने का काल निकालिए।
  \item कॉर्टिसोल के घुसने, बँधने और DNA तक पहुँचने का काल
        (मिनटों में) जोड़िए और उन देखे गए \qty{4}{h} को समझाइए।
  \item देह कॉर्टिसोल के काम के लिए ऐड्रिनैलिन जैसा कोई सोपान क्यों नहीं
        लेती, और किसी डर के लिए कॉर्टिसोल की क्रियाविधि क्यों नहीं?
  \item कॉर्टिसोल यकृत-कोशिकाओं में ऐड्रिनैलिन-ग्राहियों की संख्या भी
        चढ़ा देता है। समझाइए कि कोई धीमा \omterm{def:b2:cell-signalling:kinds}{हॉर्मोन} किसी तेज़ \omterm{def:b2:cell-signalling:kinds}{हॉर्मोन} को कैसे
        प्रवर्धित कर सकता है, और यह उस तेज़ के अधिग्रहण-वक्र का क्या करता है।
\end{enumerate}

\medskip
\noindent\textbf{भाग IV --- वह पाश।}
\begin{enumerate}[resume]
  \item किसी भोजन के बाद ग्लूकोज़ 5 से चढ़कर \qty{8}{mmol/L} हो जाती है।
        ग्लूकोज़ से इंसुलिन-मोचन तक $\beta$ कोशिका की अनुक्रिया का
        वर्णन कीजिए, और उस युग्मन का नाम लीजिए।
  \item इंसुलिन किसी टाइरोसीन काइनेज़ से होकर काम करता है, ग्लूकागॉन
        cAMP से होकर। समझाइए कि एक ही कोशिका दोनों को कैसे मान सकती है और
        यकृत के ग्लाइकोजन-एंजाइम उनके अनुपात से कैसे बैठते हैं।
  \item किसी रोगी के पास $\beta$ कोशिकाएँ नहीं हैं। उपवास तथा
        भोजन के बाद की ग्लूकोज़ की भविष्यवाणी कीजिए, और समझाइए कि वसा क्यों जलती है और अम्ल
        क्यों प्रकट होता है।
  \item इंजेक्ट की गई इंसुलिन में कोई प्रतिपुष्टि नहीं है। सामान्य ख़ुराक के साथ
        किसी असामान्य रूप से लंबी दौड़ का ख़तरा समझाइए।
  \item \cref{ch:b2:blood-pressure} के \omterm{def:b2:blood-pressure:baroreflex}{दाब-प्रतिवर्त} से
        ग्लूकोज़-पाश की तुलना कीजिए: संवेदक, प्रभावक, लब्धि, गति।
  \item परिणाम बताइए: छोड़े जाने के बाद ऐड्रिनैलिन का अधिग्रहण, वह
        सोपान-लब्धि, यकृत का ग्लूकोज़-निर्गम, और कॉर्टिसोल के लिए गया
        काल।
\end{enumerate}
\end{problem}
